\chapter{1994年上海卷}
\section{填空题}

\begin{problem}
不等式 \(|x + 1| < 1\) 的解是\fillinblank{}.

\end{problem}

\begin{answer}
\(-2 < x < 0\)
\end{answer}

\begin{problem}
计算： \(\sin \left(\frac{1}{2} \operatorname{arc} \cos \frac{1}{8}\right) =\)\fillinblank{}.

\end{problem}

\begin{answer}
\(\frac{\sqrt{7}}{4}\)
\end{answer}

\begin{problem}
以点 \(C(-2,3)\) 为圆心且与 \(y\) 轴相切的圆的方程是\fillinblank{}.

\end{problem}

\begin{answer}
\((x + 2)^2 + (y - 3)^2 = 4\)
\end{answer}

\begin{problem}
方程 \(\log_3(x - 1) = \log_9(x + 5)\) 的解是\fillinblank{}.

\end{problem}

\begin{answer}
\(x = 4\)
\end{answer}

\begin{problem}
有一个实心圆锥体的零部件，它的轴截面是边长为 \(10\) 厘米的等边三角形. 现要在它的整个表面镀上一层防腐材料，已知每平方厘米的工料价为 \(0.10\) 元，则需要费用\fillinblank{} 元（\(\pi\) 取 \(3.2\)）.

\end{problem}

\begin{answer}
\(24\)
\end{answer}

\begin{problem}
函数 \(y = \sqrt{x^2 - 1} (x \leqslant -1) \) 的反函数是\fillinblank{}.

\end{problem}

\begin{answer}
\(y = -\sqrt{x^2 + 1}\)（\(x \geqslant 0\)）
\end{answer}

\begin{problem}
双曲线 \(\frac{y^2}{2} - x^2 = 1\) 的两个焦点的坐标是\fillinblank{}.

\end{problem}

\begin{answer}
\((0,\sqrt{3})\) 和 \((0,-\sqrt{3})\)
\end{answer}

\begin{problem}
\(\left(x^{3} - \frac{2}{x^{2}}\right)^{5}\) 的展开式中 \(x^{5}\) 的系数是 （结果用数值表示）\fillinblank{}.

\end{problem}

\begin{answer}
\(40\)
\end{answer}

\begin{problem}
函数 \(y = \sin 2x - 2\cos^2 x \) 的最大值是\fillinblank{}.

\end{problem}

\begin{answer}
\(\sqrt{2} - 1\)
\end{answer}

\begin{problem}
函数 \(y = \frac{\sin 2x + \sin\left(2x + \frac{\pi}{3}\right)}{\cos 2x + \cos\left(2x + \frac{\pi}{3}\right)} \) 的最小正周期是\fillinblank{}.

\end{problem}

\begin{answer}
\(\frac{\pi}{2}\)
\end{answer}

\section{单选题}

\begin{problem}
当 \(a > 1\) 时，函数 \(y = \log_a x\) 和 \(y = (1 - a)x\) 的图象只可能是 （\quad）

\choices
    {\choicebitmap[width=0.15\paperwidth]{img/1994/shanghai/q11_optA_nolabel.png}}
    {\choicebitmap[width=0.15\paperwidth]{img/1994/shanghai/q11_optB_nolabel.png}}
    {\choicebitmap[width=0.15\paperwidth]{img/1994/shanghai/q11_optC_nolabel.png}}
    {\choicebitmap[width=0.15\paperwidth]{img/1994/shanghai/q11_optD_nolabel.png}}

\end{problem}

\begin{problem}
如果 \(0 < a < 1\)，那么下列不等式中正确的是 （\quad）

\choices
    {\((1 - a)^{\frac{1}{3}} > (1 - a)^{\frac{1}{2}}\)}
    {\(\log_{(1 - a)}(1 + a) > 0\)}
    {\((1 - a)^{3} > (1 + a)^{2}\)}
    {\((1 - a)^{1 + a} > 1\).}

\end{problem}

\begin{problem}
已知点 \(P\) 的极坐标为 \((1, \pi)\)，那么过点 \(P\) 且垂直于极轴的直线的极坐标方程为（\quad）

\choices
    {\(\rho = 1\)}
    {\(\rho = \cos \theta\)}
    {\(\rho = -\frac{1}{\cos\theta}\)}
    {\(\rho = \frac{1}{\cos\theta}\).}

\end{problem}

\begin{problem}
已知 \(a\)，\(b\) 是异面直线，直线 \(c\) 平行于直线 \(a\)，那么 \(c\) 与 \(b\) （\quad）

\choices
    {一定是异面直线}
    {一定是相交直线}
    {不可能是平行直线；}
    {不可能是相交直线.}

\end{problem}

\begin{problem}
设 \(I\) 是全集，集合 \(P\)，\(Q\) 满足 \(P \subset Q\)，则下面的结论中错误的是 （\quad）

\choices
    {\(P \cup Q = Q\)}
    {\(\overline{P} \cup Q = I\)}
    {\(P \cap \overline{Q} = \varnothing\)}
    {\(\overline{P} \cap \overline{Q} = \overline{P}\).}

\end{problem}

\begin{problem}
设复数 \(z = -\frac{1}{2} + \frac{\sqrt{3}}{2}\i\)（\(\i\) 为虚数单位），则满足等式 \(z^n = z\)，且大于 \(1\) 的正整数 \(n\) 中最小的是（\quad）

\choices
    {3}
    {4}
    {6}
    {7.}

\end{problem}

\begin{problem}
设 \(a\)，\(b\) 是平面 \(\alpha\) 外的任意两条线段，则“\(a\)，\(b\) 的长相等”是“\(a\)，\(b\) 在平面 \(\alpha\) 内的射影长相等”的 （\quad）

\choices
    {非充分条件也非必要条件}
    {充分必要条件}
    {必要条件而非充分条件}
    {充分条件而非必要条件.}

\end{problem}

\begin{problem}
计划在某画廊展出10幅不同的画，其中1幅水彩画，4幅油画，5幅国画，排成一行陈列，要求同一品种的画必须连在一起，并且水彩画不放在两端，那么不同陈列方式有（\quad）

\choices
    {\(\mathrm A_4^4 \cdot \mathrm A_5^5\) 种}
    {\(\mathrm A_3^3 \cdot \mathrm A_4^4 \cdot \mathrm A_5^5\) 种}
    {\(\mathrm C_3^1 \cdot \mathrm A_4^4 \cdot \mathrm A_5^5\) 种}
    {\(\mathrm A_2^2 \cdot \mathrm A_4^4 \cdot \mathrm A_5^5\) 种.}

\end{problem}

\begin{problem}
在直角坐标系 \(xOy\) 中，曲线 \(C\) 的方程是 \(y = \cos x\). 现平移坐标系，把原点移到点 \(O'\left(\frac{\pi}{2}, -\frac{\pi}{2}\right)\)，则在坐标系 \(x'O'y'\) 中，曲线 \(C\) 的方程是 （\quad）

\choices
    {\(y' = \sin x' + \frac{\pi}{2} \)}
    {\(y' = -\sin x' + \frac{\pi}{2} \)}
    {\(y' = \sin x' - \frac{\pi}{2} \)}
    {\(y' = -\sin x' - \frac{\pi}{2} \).}

\end{problem}

\begin{problem}
某个命题与自然数 \(n\) 有关. 如果当 \(n = k\)（\(k \in \mathbb{N}\)）时该命题成立，那么可推得当 \(n = k + 1\) 时该命题也成立. 现已知当 \(n = 5\) 时该命题不成立，那么可推得 （\quad）

\choices
    {当 \(n = 6\) 时该命题不成立}
    {当 \(n = 6\) 时该命题成立}
    {当 \(n = 4\) 时该命题不成立}
    {当 \(n = 4\) 时该命题成立.}

\end{problem}

\section{解答题}

\begin{problem}
已知 \(\sin \alpha = \frac{3}{5}\)，\(\alpha \in \left(\frac{\pi}{2}, \pi\right)\)，\(\tan(\pi - \beta) = \frac{1}{2}\)，求 \(\tan(\alpha - 2\beta)\) 的值.

\end{problem}

\begin{solution}
因为 \(\sin \alpha = \frac{3}{5}\)，\(\alpha \in \left(\frac{\pi}{2}, \pi\right)\)，所以 \(\cos \alpha = -\frac{4}{5}\)，\(\tan \alpha = -\frac{3}{4}\). 因为 \(\tan(\pi - \beta) = \frac{1}{2}\)，所以 \(\tan \beta = -\frac{1}{2}\)，从而
\[
\tan 2\beta = \frac{2\tan\beta}{1 - \tan^2\beta} = \frac{2 \cdot \left(-\frac{1}{2}\right)}{1 - \left(-\frac{1}{2}\right)^2} = -\frac{4}{3}
\]
于是
\[
\tan(\alpha - 2\beta) = \frac{\tan\alpha - \tan2\beta}{1 + \tan\alpha \tan2\beta} = \frac{-\frac{3}{4} + \frac{4}{3}}{1 + \left(-\frac{3}{4}\right)\left(-\frac{4}{3}\right)} = \frac{7}{24}
\]
\end{solution}

\begin{problem}
设 \(w\) 为复数. 它的辐角的主值为 \(\frac{3\pi}{4}\)，且 \(\frac{\left(\overline{w}\right)^2 - 4}{w}\) 为实数，求复数 \(w\).

\end{problem}

\begin{solution}
【法一】设 \(w = a + b\i\)（\(a\)，\(b \in \R\)）. 由 \(\arg w = \frac{3\pi}{4}\)，可得 \(b = -a\) 且 \(a < 0\). 又 \[\frac{\left(\overline{w}\right)^2 - 4}{w} = \frac{(a - b\i)^2 - 4}{a + b\i} = \frac{a^2(1 + \i)^2 - 4}{a(1 - \i)} = -\frac{a^2 + 2}{a} + \frac{a^2 - 2}{a}\i \] 由 \(\frac{\left(\overline{w}\right)^2 - 4}{w}\) 为实数可知 \(a^2 - 2 = 0\)，即 \(a = \pm \sqrt{2}\). 因为 \(a < 0\)，所以 \(a = -\sqrt{2}\). 于是 \(w = -\sqrt{2} + \sqrt{2}\i\).

【法二】由题意，可设 \(w = r\left(\cos \frac{3\pi}{4} + \i \sin \frac{3\pi}{4}\right)\). 则
\[
\frac{\left(\overline{w}\right)^2 - 4}{w} = \frac{\left(\overline{w}\right)^3 - 4w\overline{w}}{w\overline{w}} = \frac{r^3 \left(\cos \frac{3\pi}{4} - \i \sin \frac{3\pi}{4}\right)^3}{r^2} - \frac{4r\left(\cos \frac{3\pi}{4} - \i \sin \frac{3\pi}{4}\right)}{r^2} = \frac{(r^2 + 4) + (4 - r^2)\i}{\sqrt{2}r}
\]
由 \(\frac{\left(\overline{w}\right)^2 - 4}{w}\) 为实数，可得 \(4 - r^2 = 0\). 因此 \(r = 2\). 于是 \(w = 2\left(\cos \frac{3\pi}{4} + \i \sin \frac{3\pi}{4}\right) = -\sqrt{2} + \sqrt{2}\i.\)
\end{solution}

\begin{problem}
如图，在梯形 \(ABCD\) 中，\(AD \parallel BC\)，\(\angle ABC = \frac{\pi}{2}\)，\(AB = a\)，\(AD = 3a\)，且 \(\angle ADC = \arcsin \frac{\sqrt{5}}{5}\)，又 \(PA \perp\) 平面 \(ABCD\)，\(PA = a\). 求：
\begin{enumerate}
\item 二面角 \(P-CD-A\) 的大小（用反三角函数表示）；
\item 点 \(A\) 到平面 \(PBC\) 的距离.
\end{enumerate}

\bitmapfigure[max width=0.58\linewidth]{img/1994/shanghai/q23_fig1.png}

\end{problem}

\begin{solution}
（1）在平面 \(ABCD\) 内，过 \(A\) 作 \(AE \perp CD\)，垂足为 \(E\)，连接 \(PE\). 因为 \(PA \perp\) 平面 \(ABCD\)，由三垂线定理知 \(PE \perp CD\). 因此 \(\angle PEA\) 是二面角 \(P-CD-A\) 的平面角. 在直角三角形 \(AED\) 中，\(AD = 3a\)，\(\angle ADE = \arcsin \frac{\sqrt{5}}{5}\). 因此 \(AE = AD \cdot \sin \angle ADE = \frac{3\sqrt{5}}{5}a\). 在直角三角形 \(PAE\) 中，\(\tan \angle PEA = \frac{PA}{AE} = \frac{\sqrt{5}}{3}\)，所以 \(\angle PEA = \arctan \frac{\sqrt{5}}{3}\). 故二面角 \(P-CD-A\) 的大小为 \(\arctan \frac{\sqrt{5}}{3}\).

（2）在平面 \(PAB\) 中，过 \(A\) 作 \(AH \perp PB\)，垂足为 \(H\). 因为 \(PA \perp\) 平面 \(ABCD\)，所以 \(PA \perp BC\). 又 \(AB \perp BC\)，因此 \(BC \perp\) 平面 \(PAB\)，从而 \(BC \perp AH\). 故 \(AH \perp\) 平面 \(PBC\)，所以 \(AH\) 的长即为点 \(A\) 到平面 \(PBC\) 的距离. 在等腰直角三角形 \(PAB\) 中，\(AH = \frac{\sqrt{2}}{2}a\). 所以点 \(A\) 到平面 \(PBC\) 的距离为 \(\frac{\sqrt{2}}{2}a\).
\end{solution}

\begin{problem}
设椭圆的中心为原点 \(O\)，一个焦点为 \(F(0,1)\)，长轴和短轴的长度之比为 \(t\).
\begin{enumerate}
\item 求椭圆的方程；
\item 设经过原点且斜率为 \(t\) 的直线与椭圆在 \(y\) 轴右边部分的交点为 \(Q\)，点 \(P\) 在该直线上，且 \(\left|\frac{OP}{OQ}\right| = t\sqrt{t^2 - 1}\)，当 \(t\) 变化时，求点 \(P\) 的轨迹方程，并说明轨迹是什么图形.
\end{enumerate}

\end{problem}

\begin{solution}
（1）设所求椭圆方程为 \(\frac{x^2}{b^2} + \frac{y^2}{a^2} = 1\)（\(a > b > 0\)）. 由题意得
\[
\begin{cases}
\sqrt{a^2 - b^2} = 1, \\
\frac{a}{b} = t.
\end{cases}
\]
解得
\[
\begin{cases}
a^2 = \frac{t^2}{t^2 - 1}, \\
b^2 = \frac{1}{t^2 - 1}.
\end{cases}
\]
所以椭圆的方程为 \(t^2(t^2 - 1)x^2 + (t^2 - 1)y^2 = t^2\).

（2）设点 \(P\) 和 \(Q\) 的坐标分别为 \((x,y)\) 和 \((x',y')\). 解方程组
\[
\begin{cases}
t^2(t^2 - 1)x'^2 + (t^2 - 1)y'^2 = t^2, \\
y' = tx',
\end{cases}
\]
得
\[
\begin{cases}
x' = \frac{1}{\sqrt{2(t^2 - 1)}}, \\
y' = \frac{t}{\sqrt{2(t^2 - 1)}}.
\end{cases}
\]
由 \(\frac{|OP|}{|OQ|} = t\sqrt{t^2 - 1}\) 和 \(\frac{|OP|}{|OQ|} = \left|\frac{x}{x'}\right|\) 得
\[
\begin{cases}
x = \frac{t}{\sqrt{2}}, \\
y = \frac{t^2}{\sqrt{2}},
\end{cases}
\text{或}
\begin{cases}
x = -\frac{t}{\sqrt{2}}, \\
y = -\frac{t^2}{\sqrt{2}}.
\end{cases}
\]
而 \(t > 1\)，于是点 \(P\) 的轨迹方程为 \(x^2 = \frac{\sqrt{2}}{2}y \left(x > \frac{\sqrt{2}}{2}\right)\) 和 \(x^2 = -\frac{\sqrt{2}}{2}y \left(x < -\frac{\sqrt{2}}{2}\right)\).
点 \(P\) 的轨迹为抛物线 \(x^2 = \frac{\sqrt{2}}{2}y\) 在直线 \(x = \frac{\sqrt{2}}{2}\) 右侧的部分和抛物线 \(x^2 = -\frac{\sqrt{2}}{2}y\) 在直线 \(x = -\frac{\sqrt{2}}{2}\) 左侧的部分.
\end{solution}

\begin{problem}
在直角坐标系 \(xOy\) 中，设矩形 \(OPQR\) 的顶点按逆时针顺序依次为 \(O(0,0)\)，\(P(1,t)\)，\(Q(1-2t,2+t)\)，\(R(-2t,2)\)，其中 \(t \in (0,+\infty)\).
\begin{enumerate}
\item 求矩形 \(OPQR\) 在第一象限部分的面积 \(S(t)\)；
\item 确定函数 \(S(t)\) 的单调区间，并加以证明.
\end{enumerate}

\end{problem}

\begin{solution}
（1）当 \(-2t + 1 > 0\)，即 \(0 < t < \frac{1}{2}\) 时，如图，点 \(Q\) 在第一象限. 此时 \(S(t)\) 为四边形 \(OPQK\) 的面积，直线 \(QR\) 的方程为 \(y - 2 = t(x + 2t)\).

\bitmapfigure[max width=0.64\linewidth]{img/1994/shanghai/q25_solution_fig1.png}

令 \(x = 0\)，得 \(y = 2t^2 + 2\)，点 \(K\) 的坐标为 \((0,2t^2 + 2)\). 因而
\examdisplaychain{
S_{OPQK}
&= S_{OPQR} - S_{OKR} \\
&= 2\left(\sqrt{1 + t^2}\right)^2 - \frac{1}{2}(2t^2 + 2)\cdot 2t \\
&= 2(1 - t + t^2 - t^3).
}

当 \(-2t + 1 \leqslant 0\)，即 \(t \geqslant \frac{1}{2}\) 时，如图，点 \(Q\) 在 \(y\) 轴上或第二象限. 此时 \(S(t)\) 为 \(\triangle OPL\) 的面积，直线 \(PQ\) 的方程为 \(y - t = -\frac{1}{t}(x - 1)\). 令 \(x = 0\)，得 \(y = t + \frac{1}{t}\)，点 \(L\) 的坐标为 \(\left(0, t + \frac{1}{t}\right)\). 于是
\(S_{OPL} = \frac{1}{2}\left(t + \frac{1}{t}\right)\).

所以
\[
S(t) =
\begin{cases}
2(1 - t + t^2 - t^3), & 0 < t < \frac{1}{2}, \\
\frac{1}{2}\left(t + \frac{1}{t}\right), & t\geqslant \frac{1}{2}.
\end{cases}
\]

（2）当 \(0 < t < \frac{1}{2}\) 时，对于任何 \(0 < t_1 < t_2 < \frac{1}{2}\)，有
\(S(t_1) - S(t_2) = 2(t_2 - t_1)\left[1 - (t_1 + t_2) + \left(t_1^2 + t_1 t_2 + t_2^2\right)\right] > 0\).
即 \(S(t_1) > S(t_2)\)，所以 \(S(t)\) 在区间 \(\left(0,\frac{1}{2}\right)\) 内是减函数.

当 \(t \geqslant \frac{1}{2}\) 时，对于任何 \(\frac{1}{2} \leqslant t_1 < t_2\)，有
\(S(t_1) - S(t_2) = \frac{1}{2}(t_1 - t_2)\left(1 - \frac{1}{t_1 t_2}\right)\).
所以，若 \(\frac{1}{2} \leqslant t_1 < t_2 \leqslant 1\)，则 \(S(t_1) > S(t_2)\)；若 \(1 \leqslant t_1 < t_2\)，则 \(S(t_1) < S(t_2)\). 所以 \(S(t)\) 在区间 \(\left[\frac{1}{2},1\right]\) 上是减函数，在区间 \([1,+\infty)\) 内是增函数.

又 \(2\left[1 - \frac{1}{2} + \left(\frac{1}{2}\right)^2 - \left(\frac{1}{2}\right)^3\right] = \frac{5}{4} = S\left(\frac{1}{2}\right)\).
由上面的证明过程可得，对于任何 \(0 < t_1 < \frac{1}{2} \leqslant t_2 < 1\)，均有 \(S(t_1) < \frac{5}{4} \leqslant S(t_2)\). 于是 \(S(t)\) 的单调区间分别为 \(\left(0,1\right]\) 及 \(\left[1,+\infty\right)\)，且 \(S(t)\) 在 \(\left(0,1\right]\) 内是减函数，在 \(\left[1,+\infty\right)\) 内是增函数.
\end{solution}

\begin{problem}
已知数列 \(\{a_n\}\) 满足条件：\(a_1 = 1\)，\(a_2 = r\)（\(r > 0\)），且 \(\{a_n a_{n+1}\}\) 是公比为 \(q\)（\(q > 0\)）的等比数列. 设 \(b_n = a_{2n-1} + a_{2n}\)（\(n = 1\)，\(2\)，\(\dots\)）.
\begin{enumerate}
\item 求出使不等式 \(a_n a_{n+1} + a_{n+1} a_{n+2} > a_{n+2} a_{n+3}\)（\(n \in \mathbb{N}\)）成立时 \(q\) 的取值范围；
\item 求 \(b_n\) 和 \(\lim\limits_{n \to \infty} \frac{1}{S_n}\)，其中 \(S_n = b_1 + b_2 + \cdots + b_n\)；
\item 设 \(r = 2^{19.2} - 1\)，\(q = \frac{1}{2}\)，求数列 \(\left\{\frac{\log_2 b_{n+1}}{\log_2 b_n}\right\}\) 的最大项和最小项的值.
\end{enumerate}

\end{problem}

\begin{solution}
（1）由题意得 \(rq^{n-1} + rq^n > rq^{n+1}\). 由题设 \(r > 0\)，\(q > 0\)，故从上式可得 \(q^2 - q - 1 < 0\). 解得 \(\frac{1 - \sqrt{5}}{2} < q < \frac{1 + \sqrt{5}}{2}\). 又因 \(q > 0\)，故 \(0 < q < \frac{1 + \sqrt{5}}{2}\).

（2）因为 \(\frac{a_{n+1} a_{n+2}}{a_n a_{n+1}} = \frac{a_{n+2}}{a_n} = q\)，所以
\examdisplaychain{
\frac{b_{n+1}}{b_n}
&= \frac{a_{2n+1} + a_{2n+2}}{a_{2n-1} + a_{2n}} \\
&= \frac{a_{2n-1}q + a_{2n}q}{a_{2n-1} + a_{2n}} \\
&= q \neq 0.
}
\(b_1 = 1 + r \neq 0\)，所以 \(\{b_n\}\) 是首项为 \(1 + r\)，公比为 \(q\) 的等比数列，从而 \(b_n = (1 + r)q^{n-1}\).

当 \(q = 1\) 时，\(S_n = n(1 + r)\)，所以 \(\lim\limits_{n \to \infty} \frac{1}{S_n} = \lim\limits_{n \to \infty} \frac{1}{n(1 + r)} = 0\).

当 \(0 < q < 1\) 时，\(S_n = \frac{(1 + r)(1 - q^n)}{1 - q}\)，所以 \(\lim\limits_{n \to \infty} \frac{1}{S_n} = \lim\limits_{n \to \infty} \frac{1 - q}{(1 + r)(1 - q^n)} = \frac{1 - q}{1 + r}\).

当 \(q > 1\) 时，\(S_n = \frac{(1 + r)(1 - q^n)}{1 - q}\)，所以 \(\lim\limits_{n \to \infty} \frac{1}{S_n} = \lim\limits_{n \to \infty} \frac{1}{(1 + r)(1 - q^n)} = 0\).

所以
\[
\lim\limits_{n \to \infty} \frac{1}{S_n}
=
\begin{cases}
\frac{1 - q}{1 + r}, & 0 < q < 1, \\
0, & q \geqslant 1.
\end{cases}
\]

（3）由（2）得 \(b_n = (1 + r)q^{n-1}\)，所以
\examdisplaychain{
\frac{\log_2 b_{n+1}}{\log_2 b_n}
&= \frac{\log_2\left[(1 + r)q^n\right]}{\log_2\left[(1 + r)q^{n-1}\right]} \\
&= \frac{\log_2(1 + r) + n\log_2 q}{\log_2(1 + r) + (n - 1)\log_2 q} \\
&= 1 + \frac{1}{n - 20.2}.
}
记 \(c_n = \frac{\log_2 b_{n+1}}{\log_2 b_n}\). 从上式可知，当 \(n - 20.2 > 0\)，即 \(n \geqslant 21\)（\(n \in \mathbb{N}\)）时，\(c_n\) 随 \(n\) 的增大而减小，故
\(1 < c_n \leqslant c_{21} = 1 + \frac{1}{21 - 20.2} = 1 + \frac{1}{0.8} = 2.25\).

当 \(n - 20.2 < 0\)，即 \(n \leqslant 20\)（\(n \in \mathbb{N}\)）时，\(c_n\) 也随 \(n\) 的增大而减小，故
\(1 > c_n \geqslant c_{20} = 1 + \frac{1}{20 - 20.2} = 1 - \frac{1}{0.2} = -4\).

综合可知，对任意自然数 \(n\)，都有 \(c_{20} \leqslant c_n \leqslant c_{21}\). 故 \(\{c_n\}\) 的最大项为 \(c_{21} = 2.25\)，最小项为 \(c_{20} = -4\).
\end{solution}
